\titlespacing*{\subsection}{0pt}{1.0em}{0.3em}
\subsection{외환 환율 예측 모델의 한계점}

환율 예측 분야는 2000년대 이전부터 계량경제학적 접근을 통해 활발히 연구되었다. 초기 연구들은 자기회귀 누적 이동평균(Autoregressive Integrated Moving Average, ARIMA) 모형, 일반화 자기회귀 조건부 이분산성(Generalized Autoregressive Conditional Heteroskedasticity, GARCH)모형 과 같은 전통적인 시계열 모델들을 사용하여 환율의 시간적 구조와 변동성을 설명하고자 하였다~\cite{Bollerslev1986GARCH,Engle1982ARCH}. 그러나 이러한 방식은 통계적 가정을 필요로 하며, 실제 금융 시장의 복잡한 비선형 패턴을 포착하는 데 한계가 있다~\cite{Fischer2018}. 이에 따라서 최근에는 장단기 메모리(Long Short-Term Memory, LSTM), 합성곱 신경망(Convolutional Neural Network, CNN), Transformer, Informer 등 딥러닝 기반 방법론이 환율 예측 연구에서 폭넓게 쓰이고 있다. 국내에서도 주가 지수 예측과 모바일 결제 수요 예측을 대상으로 한 신경망 기반 시계열 예측 연구가 보고되었다~\cite{Ha2019WaveletStock,Kim2018MobileDemand}. 그러나 기존 모델을 활용한 연구들은 환율을 주로 개별 시계열로 취급하거나 입력 변수를 단순 벡터 형태로 처리함으로써, 외환 시장에서 나타나는 국가 간 상호의존적 관계 구조를 충분히 반영하지 못하는 한계가 있다. 또한 한 국가의 경제 상황이나 금융 시장 변화가 다른 국가의 환율에 연쇄적으로 영향을 미치는 전이 구조를 고정된 것으로 가정하는 경우도 존재한다. 이러한 접근 방식은 시장 환경 변화에 따라 변하는 국가 간 상호의존적 관계를 충분히 반영하지 못하며, 결과적으로 환율 변동의 복잡한 상호작용 구조를 설명하거나 장기 예측에서 안정성을 확보하는 데 한계를 가진다~\cite{Isufaj2025}.

\subsection{그래프 신경망 기반의 다변량 시계열 예측}
앞서 언급한 기존 환율 예측의 한계점을 극복하기 위해, 그래프 신경망(Graph Neural Network, GNN)을 활용한 시계열 예측 연구가 주목받고 있다~\cite{Wu2020MTGNN,Cao2020StemGNN}. 국내의 경우 ETF 가격 방향성 예측을 위한 그래프 구조 학습 기반 신경망과 신용카드 거래 데이터를 활용한 BiLSTM-GAT 기반 매출 예측 모델이 제안되며, 그래프 기반 예측의 적용 가능성이 확장되고 있다~\cite{Jo2024ETFGraph,Jung2024BiLSTMGAT}. 이 모델은 다변량 시계열 데이터에서 각 변수를 그래프의 노드로 간주하여 노드 간의 공간적 의존성을 효과적으로 학습한다~\cite{Kipf2017GCN,Wu2019GraphWaveNet}. 특히 변수 간 상호의존성을 명시적으로 모델링함으로써 기존의 CNN 기반 방법론보다 높은 해석력을 제공한다~\cite{Velickovic2018GAT}. 다만 외환 시장과 같은 일반적인 다변량 시계열 데이터는 변수 간의 연결 구조가 사전에 정의되어 있지 않다. 이를 해결하기 위해 Wu et al.(2020)은 사전 정의된 그래프 구조 없이 데이터로부터 변수 간의 관계를 자동으로 추출하는 MTGNN 프레임워크를 제안하였다~\cite{Wu2020MTGNN}. 그러나 MTGNN은 기본적으로 단일 예측값을 출력하는 결정론적 모델로, 예측 결과에 대한 불확실성을 직접적으로 제공하지 않는다. 이러한 한계는 구조적 변동성과 높은 노이즈를 포함하는 외환 시장과 같은 금융 시계열 데이터에서 더 두드러진다. 금융 데이터에서는 예측값뿐 아니라 예측의 신뢰도를 함께 고려할 필요가 있기 때문이다. 이러한 한계를 보완하기 위해 본 연구에서는 MTGNN 구조에 베이지안 접근을 결합한 B--MTGNN(Bayesian Multivariate Time-series Graph Neural Network)을 제안한다.  

\subsection{멀티 에이전트 시스템}
최근 그래프 기반 예측 연구는 불확실성까지 포함한 정량 예측이 가능해졌지만, 경제적 맥락에서 그 결과물을 해석하기는 여전히 어렵다는 문제가 있다. 국내 금융 예측 연구에서도 설명 가능성을 높이기 위해 LRP 기반 KOSPI 증감 예측 모델이 제안된 바 있다~\cite{Lee2021KOSPILRP}. 따라 최근 연구들은 LLM을 활용해 정량적 모델의 맥락 이해 한계를 보완하려는 시도가 꾸준히 이어지고 있다. 그러나 환율과 같이 외부 요인(경제 리스크, 통화정책 등)에 민감하게 작용하는 도메인에서, 단일 LLM에 의존하는 방식은 고도의 전문 지식을 종합하는 과정에서 대표적인 문제로 지적되고 있는 환각이나 확증 편향에 빠지기 쉽다는 치명적인 한계를 드러냈다.
이러한 단일 모델의 제약을 극복하기 위해 의료(ClinicalRAG), 법률(L-MARS), 보안(AutoBnB-RAG) 등 오류에 민감한 분야를 중심으로 전문화된 멀티 에이전트 시스템이 연구되고 있다. 의료 분야에서는 ClinicalRAG~\cite{ClinicalRAG2024}와 같이 다중 에이전트 구조를 통해 방대한 임상 데이터를 교차 검증하여 임상 의사결정 지원을 통해 진단의 신뢰성을 높이는 시도가 성공적으로 이루어졌으며, 법률 도메인에서는 L-MARS~\cite{Wang2025LMARS}와 같은 시스템이 복잡한 판례 해석 작업에서 에이전트 간 역할 분리를 통해 추론의 일관성을 확보하고 환각 현상을 획기적으로 줄였다. 또한 사이버 보안 분야의 AutoBnB-RAG~\cite{Liu2025AutoBnBRAG}는 다중 에이전트 기반의 동적 검색을 통해 지능형 위협에 대한 사후 분석 역량을 확인하였다. 이들 연구는 각자 전문성을 가진 에이전트들이 각자의 전문 지식 도메인에서 정보를 처리함으로써 단일 모델이 감당하기 어려운 인지적 부담을 해소하고 추론의 심도와 일관성을 비약적으로 향상할 수 있음을 입증하였다~\cite{Wu2023AutoGen}. 
그럼에도, 정밀한 거시경제 지표를 활용하여 정량적 예측 결과를 사후적으로 평가하고 해석하는 멀티 에이전트 시스템을 적용한 사례는 여전히 부족하다. 기존의 멀티 에이전트 연구들은 처리량 증가, 결과 도출 속도 저하, 에이전트 간 동조현상으로 인한 오류 전파 및 노이즈 증폭의 치명적인 단점을 초래하였다~\cite{Cui2025FreeMAD}. 결과적으로 현재의 외환 환율 예측은 고도화된 학습 모델을 통한 정량적 패턴 인식과, 전문화된 시스템을 통한 정성적 검증을 결합해야 하는 단계에 이르렀다. 본 연구는 이러한 기존 방법론들의 발전 과정과 한계를 종합하여, 딥러닝 기반의 다변량 시계열 예측 결과에 전문화된 멀티 에이전트 시스템을 결합하는 방법론을 제안한다. 이를 통해 기존 연구들이 지녔던 단일 모델의 인지적 과부하 및 오류 전파 문제를 극복하고, 환율 예측의 실무적 처리 효율성과 사후적 추론의 심도를 동시에 확보하고자 한다.
